% 《大学物理学学习指导》第5版第11章“变化的电磁场”，例11.1—11.7。
% 原书印刷155—162页 / PDF物理163—170页，高分辨率源扫描逐例核实。
% 研究性中文短题意摘要及独立推导，不复制原始页面和图像。
\section*{变化电磁场例题：感应电动势、磁能和坡印廷矢量}

\sourceexample{11.1}{155--156}{163--164}
\problemtext 无限长竖直导线载电流$I=5$A，右侧共面的矩形线圈
高度$\ell=0.06$m、径向宽度$a=0.04$m，以$v=0.03$m/s向右匀速远离。
当近边距导线$d=0.05$m时，求感应电动势的大小和方向。
\workings
距离长直线$r$处的磁场为
$B(r)=\mu_0I/(2\pi r)$，在右侧指向纸内。
以纸内为正法向，磁通
\[
\Phi_B(x)=\ell\int_x^{x+a}\frac{\mu_0I}{2\pi r}\,dr
=\frac{\mu_0I\ell}{2\pi}\ln\frac{x+a}{x}.
\]
对$x'=v$求导，
\[
\mathcal E=-\frac{d\Phi_B}{dt}
=\frac{\mu_0I\ell av}{2\pi x(x+a)}.
\]
代入$x=d$：
\[
\boxed{|\mathcal E|=\frac{\mu_0I\ell av}{2\pi d(d+a)}
=1.60\times10^{-8}\,\mathrm V}.
\]
向纸内磁通因右移而减小，感应电流应产生纸内磁场，
从纸外观察沿\textbf{顺时针}$A\to B\to C\to D\to A$。
\crosscheck 磁通积分法和两条与长导线平行的有源边
动生电动势之差，均应给相同的结果；
若把整个线圈视为均匀磁场将丢失对数项。

\sourceexample{11.2}{157--158}{165--166}
\problemtext 半径$R$的圆柱区域内具有指向纸内、随时间均匀增大的磁感应强度，
区域外磁场为零。水平金属杆$AB$长$2R$，
其中$AC$一段长$R$位于圆内，$CB$位于圆外，
原图圆心$O$与圆内端点的几何位置构成$60^\circ$角。
求杆两端感生电动势大小及方向。
\workings
记$k=dB/dt>0$为\emph{纸内磁场大小}变化率。
以圆心为原点，杆取水平$x$轴方向，
圆内所对应的弦线可取$y=-h$，$h=\sqrt3R/2$，
内段端点$x_A=-R/2$、$x_C=R/2$。
轴对称感生电场在圆内
$\bm E_{\rm in}=(k/2)(-y\bm e_x+x\bm e_y)$，因此
\[
\mathcal E_{AC}=\int_{-R/2}^{R/2}\frac{kh}{2}\,dx
=\boxed{\frac{\sqrt3}{4}kR^2}.
\]
圆外$E_\theta=kR^2/(2r)$，沿杆$x$分量
$E_x=kR^2h/[2(x^2+h^2)]$。
从$x_C=R/2$积分到$x_B=3R/2$，
利用$\arctan(\sqrt3)-\arctan(1/\sqrt3)=\pi/6$，
\[
\mathcal E_{CB}
=\frac{kR^2}{2}\frac\pi6
=\boxed{\frac{\pi}{12}kR^2}.
\]
故全杆
\[
\boxed{\mathcal E_{AB}
=\left(\frac{\sqrt3}{4}+\frac{\pi}{12}\right)R^2\frac{dB}{dt}},
\]
感生电动势从\textbf{A指向B}，B端电势较高。
\crosscheck 圆外局部虽$B=0$，仍有非零涡旋电场，
不可将$CB$段当成零贡献。
几何角$60^\circ$在原图以半径/弦关系给出，须保留图示方向的核验依据。

\sourceexample{11.3}{158--159}{166--167}
\problemtext 电阻$R_{\mathrm e}=2\,\Omega$、面积
$S=400\,\mathrm{cm^2}$的矩形单匝线圈，
以$\omega=10\,\mathrm{rad/s}$绕垂直于
均匀磁场$B=0.5$T的轴匀速转动。
求最大磁通量、最大感应电动势、最大磁力矩，
并证明外力每周做功等于线圈的焦耳热损耗。
\workings
以法向与磁场夹角$\theta=\omega t+\theta_0$表示瞬时磁通：
$\Phi=BS\cos\theta$。
$S=400\times10^{-4}=0.04\,\mathrm{m^2}$，因此
\[
\boxed{\Phi_{\max}=BS=2.0\times10^{-2}\,\mathrm{Wb}}.
\]
感应电动势
$\mathcal E=BS\omega\sin\theta$，其峰值
\[
\boxed{\mathcal E_{\max}=BS\omega=0.20\,\mathrm V}.
\]
线圈电流大小$|I|=|\mathcal E|/R_{\mathrm e}$，
磁阻力矩的幅值为
$M=(B^2S^2\omega/R_{\mathrm e})\sin^2\theta$，
所以
\[
\boxed{M_{\max}=\frac{B^2S^2\omega}{R_{\mathrm e}}
=2.0\times10^{-3}\,\mathrm{N\,m}}.
\]
外力每转一周所做功
\[
W_{\rm ext}=\int_0^{2\pi}M\,d\theta
=\frac{\pi B^2S^2\omega}{R_{\mathrm e}}.
\]
而一周电阻中消耗的焦耳热
\[
Q=\int_0^{2\pi/\omega}I(t)^2R_{\mathrm e}\,dt
=\frac{B^2S^2\omega^2}{R_{\mathrm e}}
\cdot\frac\pi\omega=W_{\rm ext}.
\]
\crosscheck 平均$\sin^2\theta=1/2$，
净机械输入完全化为欧姆热，符合能量守恒；
需要恒定角速率的外力克服电磁阻力矩。

\sourceexample{11.4}{159--160}{167--168}
\problemtext 空气芯密绕矩形截面环形螺线管，
总匝数$N$、内半径$a$、外半径$b>a$、截面高$h$，
求自感系数，并用磁链、自感电动势与场能三种方法核对。
\workings
理想磁环内部磁场
$B(r)=\mu_0NI/(2\pi r)$，$a<r<b$，
一匝所链磁通
\[
\Phi_1=h\int_a^b B(r)dr
=\frac{\mu_0NIh}{2\pi}\ln\frac ba.
\]
总磁链$\Psi=N\Phi_1$，故
\[
\boxed{L=\frac\Psi I
=\frac{\mu_0N^2h}{2\pi}\ln\frac ba}.
\]
由感应电动势可核对
$\mathcal E_L=-d\Psi/dt=-L\,dI/dt$。
从体积元$dV=2\pi rh\,dr$积分磁能，
\[
W=\int_a^b\frac{B(r)^2}{2\mu_0}\,2\pi rh\,dr
=\frac{\mu_0N^2I^2h}{4\pi}\ln\frac ba
=\frac12LI^2.
\]
\crosscheck 三种定义一致；若磁芯磁导率$\mu$为线性常数，
把真空$\mu_0$替换为$\mu$即可，但磁滞材料一般不适用。

\sourceexample{11.5}{160--161}{168--169}
\problemtext 真空平行圆板电容器的极板半径为$R=0.10$m。
充电时板间电场的增长率为
\[
\frac{dE}{dt}=1.0\times10^{13}\,\mathrm{V/(m\,s)}.
\]
求板间位移电流，以及板间距轴线$r<R$处与$r=R$处的磁感应强度。
\workings
忽略边缘效应，$\bm D=\varepsilon_0\bm E$，
\[
\boxed{I_D=\varepsilon_0\pi R^2\frac{dE}{dt}
\approx2.78\,\mathrm A}.
\]
取板间半径$r<R$的圆形安培环路，
$\oint\bm B\cdot d\bm l=\mu_0I_{D,\rm enc}$，
$I_{D,\rm enc}=\varepsilon_0\pi r^2\,dE/dt$，因此
\[
\boxed{B(r)=\frac{\mu_0\varepsilon_0r}{2}\frac{dE}{dt}},
\qquad
\boxed{B(R)\approx5.56\times10^{-6}\,\mathrm T}.
\]
\crosscheck 电容器板间没有穿过介质的自由电子传导电流；
必须在推广安培定理中计入位移电流。

\sourceexample{11.6}{161--162}{169--170}
\problemtext 证明对几何、介电参数不随时间改变的理想平行板电容器，
板间位移电流满足$I_D=C\,dU/dt$。
\workings
固定板距$d$、面积$S$、均匀介电常数$\varepsilon$，
$E=U/d$、$D=\varepsilon E$、$C=\varepsilon S/d$，
故电位移通量
$\Phi_D=DS=\varepsilon SU/d=CU$。
由位移电流定义
\[
\boxed{I_D=\frac{d\Phi_D}{dt}
=C\frac{dU}{dt}}.
\]
\crosscheck 若$C$也随时间改变，正确求导应是
$I_D=d(CU)/dt=C\dot U+U\dot C$，
不能无条件沿用固定电容公式。

\sourceexample{11.7}{162}{170}
\problemtext 半径$a$、电阻率$\rho$的圆柱形长导体，
总直流$I_0$均匀分布于横截面。
求导体内部距轴$r$处的电场强度$E$、磁场强度$H$、
坡印廷矢量$\bm S$，并证明长度$\ell$、半径$r\le a$
的同轴圆柱部分从侧面流入的电磁功率等于其焦耳热功率。
\workings
电流密度$j=I_0/(\pi a^2)$，
欧姆定律使电场沿电流轴向，
\[
\boxed{E=\rho j=\frac{\rho I_0}{\pi a^2}}.
\]
半径$r<a$的安培环路包围
$I_{\rm enc}=I_0r^2/a^2$，
故
\[
\boxed{H(r)=\frac{I_0r}{2\pi a^2}}.
\]
电场沿导体轴线，磁场绕轴环向，因此
$\bm S=\bm E\times\bm H$
指向\textbf{径向内侧}，
\[
\boxed{|\bm S(r)|=EH
=\frac{\rho r I_0^2}{2\pi^2a^4}}.
\]
侧面进入的功率为
\[
P_{\rm in}=|\bm S(r)|\,2\pi r\ell
=\boxed{\frac{\rho I_0^2r^2\ell}{\pi a^4}}.
\]
半径$r$内部的导电电流为$I_r=I_0r^2/a^2$，
但这里所有同心柱体均承受\emph{相同轴向电场}，
其电阻$R_r=\rho\ell/(\pi r^2)$，
\[
\boxed{P_{\rm Joule}=I_r^2R_r
=\frac{\rho I_0^2r^2\ell}{\pi a^4}=P_{\rm in}}.
\]
\crosscheck 能流穿过柱体\emph{曲面侧壁}从外部进入导体，
不能把“直流电源的能量沿金属内部电子漂移”作为唯一能流解释。
